\documentclass[10pt,a4paper]{amsart}
\usepackage[margin=19mm]{geometry}
\usepackage{amsmath,amssymb,mathtools}
\usepackage[scr=rsfs,cal=euler]{mathalfa}
\usepackage{xfrac}
\usepackage[unicode,hidelinks,bookmarks=false]{hyperref}
\newcommand{\ie}{i.e.\ }
\newcommand{\cf}{cf.\ }
\newcommand{\resp}{respectively\ }
\newcommand{\Subsection}{\subsection}
\providecommand{\nobreakdash}{\nobreak\hbox{-}}
\setlength{\emergencystretch}{2em}
\multlinegap0pt
\usepackage{amssymb}
\usepackage{xcolor}
\usepackage{bm}
\usepackage{tikz}
\newtheorem{proposition}{Proposition}
\DeclareMathOperator{\Aut}{Aut}
\newcommand{\GG}{\mathcal{G}}
\DeclareMathOperator{\Sym}{\mathbb{S}}
\newcommand{\Z}{\mathbb{Z}}
\newcommand{\bb}[1]{\boldsymbol{#1}}
\title{Asymptotic number of edge-colored regular graphs}
\author{Michael Borinsky, Chiara Meroni, Maximilian Wiesmann}
\date{Source: arXiv:2601.18994v1 (CC BY 4.0)}
\begin{document}
\maketitle

\noindent\emph{Source and licence.} Asymptotic number of edge-colored regular graphs. Version arXiv:2601.18994v1 (CC BY 4.0), distributed by arXiv under the Creative Commons Attribution 4.0 International licence, http://creativecommons.org/licenses/by/4.0/, as recorded in the arXiv licence metadata for this exact version and shown on the abstract page https://arxiv.org/abs/2601.18994v1. Full licence notice: \texttt{licenses/arxiv-2601.18994-CC-BY-4.0.txt}.\par\smallskip

\noindent\textit{Editorial scope.} Counted unit: the article's proposition prop:gf\_graphs (source0.tex lines 315--325). The article's complete original proof of it is delivered with it (source0.tex lines 328--389, one \texttt{proof} environment of 62 lines). No mathematics is written, repaired or re-derived: the statement and the proof are the article's own lines, in the article's own notation, and no retained line is substituted or re-flowed.  No further source-local context is needed: every cross-reference of every retained line resolves inside the two delivered spans.  Nothing was omitted from any retained span, so the package carries no omission record.  No retained line contains a citation, so the wrapper carries no bibliography and no bibliography entry.  No retained line contains a figure environment, an image-inclusion command or drawing code, so no image is delivered and the package declares no asset dependency.  Editorial formatting only, in addition: an AMS \texttt{amsart} preamble that restates the article's own declarations and macros used by the retained lines, namely the environments \texttt{proposition} on separate counters, together with the standard packages those lines need.  The retained environments print in this document as Proposition 1 in order of appearance, whereas the article prints the same items with the numbering of its own full document.  The complete native source of the article as distributed in the arXiv e-print for this exact version is bundled unchanged as \texttt{sources/193429/source0.tex}.
\begin{proposition}
  \label{prop:gf_graphs}
The following counting formula holds:
  \begin{equation}\label{eq:gf_graphs}
    \sum_{G \in \GG} \frac{\prod_{v\in V_G} \lambda_{\deg(v)}}{|\Aut(G)|} =  
\sum_{\bb{s} \in \Z_{\geq 0}^c} 
\left(\prod_{i=1}^c (2s_i-1)!! \right)
\cdot [\bb{x}^{2\bb{s}}]\exp\left( \sum_{\bb{w}\in \Z^c_{\geq 0},\, |\bb{w}| \geq 1} \lambda_{\bb{w}} \frac{\bb{x}^{\bb{w}}}{\bb{w}!} \right),
  \end{equation}
where $\GG$ is the set of isomorphism classes of edge-colored half-edge labeled graphs. 
\end{proposition}

\begin{proof}
We interpret the right-hand side as a weighted enumeration of edge-colored graphs using the half-edge definition. Expansion of the exponential on the right-hand side and extracting the coefficient of $\bb x^{2\bb s}=x_1^{2s_1}\cdots x_c^{2s_c}$ gives
\begin{equation}\label{eq:exponential_and_block_partitions}
    [x_1^{2s_1}\cdots x_c^{2s_c}]
    \exp\left( \sum_{\bb{w}\in \Z^c_{\geq 0},\, |\bb{w}| \geq 1} \lambda_{\bb{w}} \frac{\bb{x}^{\bb{w}}}{\bb{w}!} \right)
    = 
    \sum_{\{t_{\bb{u}}\}} 
    \frac{1}{\prod_{\bb{u}} t_{\bb{u}}!\, \bb{u}!^{t_{\bb{u}}}} \prod_{\substack{\bb{u}\in \Z^c_{\geq 0}}} \lambda_{\bb{u}}^{t_{\bb{u}}},
\end{equation}
where the last sum ranges over all functions
$\bb u\mapsto t_{\bb u}\in\Z_{\ge0}$ satisfying
$\sum_{\bb u}t_{\bb u}u_i=2s_i$ for each $i=1,\dots,c$.
This condition implies that only finitely many $t_{\bb u}$ are nonzero.

Fix a tuple of sets $(H_1,\ldots,H_c)$ with $|H_i|=2s_i$ and write $H:=H_1\sqcup\cdots\sqcup H_c$.
A choice of multiplicities $\{t_{\bb u}\}$ determines a class of partitions of $H$ into blocks such that, for each $\bb u=(u_1,\dots,u_c)$, exactly $t_{\bb u}$ blocks
contain $u_i$ elements from $H_i$.
Each block represents a vertex with multidegree $\bb u$.
The number of such partitions is
\begin{equation}\label{eq:block_partitions}
\frac{(2\bb s)!}{\prod_{\bb u}t_{\bb u}!\,\bb u!^{t_{\bb u}}},
\end{equation}
which follows from the orbit--stabilizer theorem applied to the following natural action. The $c$-fold product of symmetric groups $\Sym_{2s_1} \times \dots \times \Sym_{2s_c}$ acts  on the set of partitions of $H_1 \sqcup \dots \sqcup H_c$. A partition with $t_{\bb{u}}$ blocks with $u_i$ elements from $H_i$ is stabilized by a subgroup isomorphic to $(\Sym_{u_1} \times \dots \times \Sym_{u_c})^{t_{\bb{u}}} \rtimes \Sym_{t_{\bb{u}}}$, permuting the elements inside each block and the blocks themselves.
Thus the coefficient \eqref{eq:exponential_and_block_partitions} enumerates vertex sets with
prescribed multidegrees $\bb u$ such that $\sum_{\bb u}t_{\bb u}u_i=2s_i$, weighted by the factor
$\prod_v\lambda_{\deg(v)}$.

For each color $i$, an edge of color $i$ is obtained by pairing two half-edges in
$H_i$.
The number of perfect matchings on $H_i$ is $(2s_i-1)!!=(2s_i-1) \cdot (2s_i-3) \cdots 3 \cdot 1$, and these choices are
independent for different colors.
Choosing such matchings produces a collection of colored edges; self-loops and
multiple edges are allowed, since paired half-edges may lie in the same block or in
distinct blocks.
The total number of ways to pair all half-edges into edges is therefore
\[
\prod_{i=1}^c(2s_i-1)!! \, .
\]

Combining a partition of $H$ with matchings on each $H_i$ yields a half-edge labeled
edge-colored graph which represents an isomorphism class of a graph $G$ 
with
exactly $s_i$ edges of color $i$.
For such a fixed isomorphism class, the group
$\Sym_{2s_1}\times\cdots\times\Sym_{2s_c}$
acts freely on the set of half-edge labelings by permuting half-edges within
each color.
The stabilizer of $G$ is precisely the automorphism group $\Aut(G)$.
By the orbit--stabilizer theorem, the number of half-edge labeled graphs representing $G$ is
\[
\frac{(2\bb s)!}{|\Aut(G)|}.
\]
Multiplying \eqref{eq:exponential_and_block_partitions} by the factor
\[
\frac{1}{(2\bb s)!}\prod_{i=1}^c(2s_i-1)!!
\]
has the effect of pairing half-edges into edges and then forgetting the labeling of
half-edges.
Consequently, each isomorphism class $G$ appears with weight
$\prod_{v\in V_G}\lambda_{\deg(v)}/|\Aut(G)|$.
Summing over all $\bb s\in\Z_{\ge0}^c$ proves the stated identity.
\end{proof}

\end{document}
